
\documentclass[11pt,oneside]{book}
\usepackage{../preamble}

\title{Idea Theory, Cognitive Science, and Philosophy of Mind\\
\large Volume 4 of the IdeaTheory monograph}
\author{The IdeaTheory Project}
\date{\today}

\begin{document}
\maketitle

\begin{abstract}
This fourth volume of the \emph{IdeaTheory} monograph interprets the
purely mathematical apparatus of Volume~1 --- the graded idea monoid
$(\mathcal{I},\circ,\mathbf{1})$, the bilinear resonance pairing
$\langle{-},{-}\rangle$, the Aufhebung decomposition
$x\circ y = \mathrm{sub}(x,y)\oplus \mathrm{nov}(x,y)$, the emergence
$2$-cocycle $\varepsilon$, the meaning curvature form
$\kappa(x,y)=\langle x,y\rangle-\langle y,x\rangle$, the conjugation
action, transmission chains, and structural equivalence --- as a
formal language for thought, mental content, and cognitive
architecture.  Where the social and humanistic volumes treat ideas as
public, intersubjective objects, this volume turns the same algebra
inward and asks what it says about \emph{minds that bear ideas}: how
concepts compose into thoughts, how thoughts inherit and revise prior
content, how attention and salience act as a resonance pairing, how
metacognition appears as conjugation, and how the curvature
$\kappa$ measures the irreducibly directional, non-commutative
character of inference and association documented across the cognitive
sciences.

The volume bridges a substantial non-formalized literature ---
Fodor's language of thought, Smolensky's tensor-product
representations, G\"ardenfors's conceptual spaces, Friston's
free-energy and predictive-processing program, Dennett's
intentional stance, Clark's extended mind, Lakoff and Johnson on
conceptual metaphor, Sperber's epidemiology of representations, and
the Bayesian brain tradition --- to the machine-checked theorems of
the Lean~4 development.  Each chapter follows the project-wide
template: a \emph{Background from the Literature} section grounding
the discussion in primary sources from \texttt{references.bib}, a
\emph{From Informal Claims to Formal Statements} translation table
mapping each informal cognitive thesis to a specific Lean
\texttt{theorem} or \texttt{def}, worked examples, and a
\emph{Further Reading} pointer to the relevant Volume~1 results and
their formal counterparts.

Volume~4 contains two chapters.  Chapter~7, \emph{The Algebra of
Thought}, develops the cognitive interpretation of composition,
resonance, and Aufhebung: concept combination, working-memory
binding, conceptual blending, and the asymmetry of priming.
Chapter~8, \emph{Curvature, Conjugation, and the Architecture of
Mind}, interprets the curvature form, the conjugation action, and
transmission chains as accounts of inferential asymmetry, perspective
taking, metacognition, predictive processing, and the
intergenerational transmission of mental content.  Together they
show that the same nine theorems that organise Volume~1 organise,
without modification, a surprisingly large fragment of contemporary
philosophy of mind and cognitive science --- and that every claim we
make about minds is anchored, by an explicit bridge table, to a
theorem that has been verified in Lean~4.
\end{abstract}

\tableofcontents

\chapter*{Preface to Volume 4}
\addcontentsline{toc}{chapter}{Preface to Volume 4}

\section*{Motivation}
The mathematical core of \emph{Idea Theory} was deliberately built
without reference to any particular interpretation: the nine
theorems of Volume~1 speak only of a graded monoid, a bilinear
form, a decomposition, and a cocycle.  Their power lies precisely in
their neutrality.  Volume~4 takes up the interpretive challenge that
this neutrality invites in the case of \emph{minds}.  When a thinker
entertains a thought, that thought has constituents; when those
constituents are recombined, the result is rarely the bare sum of
its parts; when one idea follows another, the order matters; when a
belief is held by an agent, it is held \emph{from a perspective} and
can be re-expressed in another; and when ideas are passed from
teacher to student, parent to child, or culture to culture, they
mutate in ways that are nevertheless lawful.  Each of these familiar
observations from cognitive science and philosophy of mind has, we
shall argue, a precise counterpart in the Volume~1 algebra.

\section*{Intended Audience and Prerequisites}
The volume is written for three overlapping audiences: philosophers
of mind looking for a formal vocabulary that does not collapse
either into pure logic or into pure neural dynamics; cognitive
scientists who want a compositional algebra compatible with both
symbolic and subsymbolic models; and formally minded readers of
Volume~1 who wish to see the abstract apparatus exercised on
concrete cognitive phenomena.  We assume familiarity with the
notation of Volume~1 (\S\S 1--3) and a working acquaintance with
either classical cognitive science (Fodor, Pylyshyn) or the
predictive-processing tradition (Clark, Friston).  No Lean
expertise is required to read the prose, but every claim is
hyperlinked, via the bridge tables, to the corresponding Lean
identifier in the namespace \texttt{IdeaTheory}.

\section*{Reading Order and Relation to the Other Volumes}
Volume~4 sits between the social/humanistic volumes (2 and 3),
which treat ideas as public objects, and the emergence volume
(5), which treats the cocycle $\varepsilon$ in its full generality.
Readers coming from Volume~1 may proceed directly to Chapter~7.
Readers coming from Volume~2 or~3 will find Chapter~7 a natural
continuation: where those volumes asked how ideas live in
\emph{communities}, this volume asks how they live in \emph{heads}.
Readers primarily interested in emergence may read Chapter~8
(especially \S 8.3 on predictive processing) before Volume~5, which
generalises the same machinery beyond the cognitive case.
Volume~6 (\emph{Advanced Applications and Unifying Perspectives})
revisits several Chapter~8 themes from a higher categorical
vantage; readers may safely defer it.

\section*{Guided Tour of the Chapters}
\textbf{Chapter 7 --- The Algebra of Thought.}  Opens with a
historical survey of compositionality debates from Frege through
Fodor and Pylyshyn's systematicity argument to Smolensky's
tensor-product response and G\"ardenfors's geometric alternative.
We then show that the monoid law $(\texttt{Theorem 1})$ formalises
the minimal commitments shared by all parties; that the resonance
pairing $(\texttt{Theorem 2})$ recovers a graded notion of semantic
similarity compatible with both vector-space and feature-based
models; and that the Aufhebung decomposition $(\texttt{Theorem 3})$
gives a precise sense to the long-standing intuition that
conceptual combination both \emph{preserves} and \emph{transcends}
its inputs --- the formal counterpart of Hampton's emergent
features and Fauconnier and Turner's conceptual blending.

\textbf{Chapter 8 --- Curvature, Conjugation, and the Architecture
of Mind.}  Takes up the four remaining Volume~1 theorems.  The
emergence cocycle $(\texttt{Theorem 4})$ and the curvature form
$(\texttt{Theorem 5})$ are interpreted as the formal trace of
order effects in reasoning (Tversky and Kahneman; Hogarth and
Einhorn) and of the directionality of associative priming.  The
conjugation action $(\texttt{Theorem 6})$ formalises perspective
taking, theory of mind, and Dennett's intentional stance; we
discuss how a self-applied conjugation yields a tractable model of
metacognition.  Transmission chains $(\texttt{Theorem 7})$ and
structural equivalence $(\texttt{Theorem 8})$ are linked to
predictive-processing accounts of belief updating and to Sperber's
epidemiology of representations.  The grading theorem
$(\texttt{Theorem 9})$ closes the volume by giving a complexity
measure on thoughts that aligns with empirical findings on
working-memory load.

\section*{Bridge to the Formalization}
Every theorem cited above lives in the Lean~4 namespace
\texttt{IdeaTheory} and is fully machine-checked.  Each chapter
ends with a \emph{From Informal Claims to Formal Statements} table
that names the Lean identifier (\texttt{IdeaTheory.composition\_assoc},
\texttt{IdeaTheory.resonance\_bilinear}, etc.) realising each
informal claim.  The appendix to this volume catalogues the
relevant Lean files.  In this way Volume~4 honours the project's
core mandate: no claim about mind is made that is not anchored in a
theorem that a computer has checked.

\input{ch7}
\input{ch8}

\appendix
\chapter{Lean 4 Files Relevant to Volume 4}
\label{app:vol4-lean}

The cognitive interpretations developed in this volume rest entirely
on the nine theorems formalised in Volume~1.  Volume~4 introduces no
new Lean files of its own; instead, every bridge table in
Chapters~7--8 cites identifiers from the files below.  Readers may
consult the source tree under \texttt{IdeaTheory/} in the project
repository.

\paragraph{\texttt{IdeaTheory/Foundations.lean}.}
Establishes the carrier type \texttt{Idea}, the composition operator
\texttt{(\(\circ\))}, the unit \texttt{1}, and the grading map
\texttt{deg : Idea \(\to\) \(\mathbb{N}\)}.  Chapter~7 cites this
file when introducing the cognitive reading of composition as
concept combination and of the unit as the ``empty thought.''

\paragraph{\texttt{IdeaTheory/Composition.lean}.}
Contains \texttt{composition\_assoc} and \texttt{composition\_unit}
(Theorem~1).  Cited in \S 7.1 to justify the systematicity of
thought in the sense of Fodor and Pylyshyn (1988).

\paragraph{\texttt{IdeaTheory/Resonance.lean}.}
Defines the bilinear resonance pairing
\texttt{resonance : Idea \(\to\) Idea \(\to\) \(\mathbb{R}\)} and
proves its bilinearity (Theorem~2).  Cited in \S 7.2 as the formal
backbone of semantic similarity and priming.

\paragraph{\texttt{IdeaTheory/Aufhebung.lean}.}
Provides \texttt{sub}, \texttt{nov}, and the decomposition theorem
\texttt{aufhebung\_decomp} (Theorem~3).  Cited in \S 7.3 as the
formalization of emergent features in conceptual combination
(Hampton, 1987) and conceptual blending (Fauconnier and Turner,
2002).

\paragraph{\texttt{IdeaTheory/Cocycle.lean}.}
Defines the emergence cocycle \texttt{epsilon} and proves the
$2$-cocycle identity (Theorem~4).  Cited in \S 8.1 as the algebraic
account of order effects in judgment.

\paragraph{\texttt{IdeaTheory/Curvature.lean}.}
Introduces \texttt{kappa x y := resonance x y - resonance y x} and
establishes its antisymmetry and exactness properties (Theorem~5).
Cited in \S 8.2 to formalise asymmetric priming and inferential
directionality.

\paragraph{\texttt{IdeaTheory/Conjugation.lean}.}
Provides the conjugation action of an idea on the algebra
(Theorem~6).  Cited in \S 8.3 as a model of perspective taking and
metacognition.

\paragraph{\texttt{IdeaTheory/Chains.lean}.}
Defines transmission chains and proves their composition law
(Theorem~7).  Cited in \S 8.4 in connection with predictive
processing and Sperber's epidemiology of representations.

\paragraph{\texttt{IdeaTheory/Equivalence.lean}.}
Defines structural equivalence \texttt{(\(\sim\))} and proves it is
a congruence with respect to composition (Theorem~8).  Cited in
\S 8.5 to model intertranslatability of cognitive idioms.

\paragraph{\texttt{IdeaTheory/Grading.lean}.}
Establishes the graded-algebra structure (Theorem~9).  Cited in
\S 8.6 as a complexity measure on thoughts compatible with
working-memory findings.

\bibliographystyle{alpha}
\bibliography{../references}

\end{document}
=== END FILE ===
